% -*- latex -*-
%
% This is ae top-level LaTeX file to be used in testing TeX building.
% It includes testing.tex which you are to provide.
% That file contains the things you want to test out, such as
% asymptote examples.
%
% Often you start out with the auto-generated file "documentation.tex"
% and then strip it down to the parts you want to scrutenize
% The results of running this test are put in mathics-test.pdf
%
% To build this:
%    make mathics-test.pdf

\documentclass[index=totoc,fleqn,10pt,oneside,toc=flat]{scrbook}

%\usepackage[utf8]{inputenc}
\usepackage[T1]{fontenc}
\usepackage{lmodern}
\usepackage[english]{babel}
\usepackage{makeidx}
\usepackage{graphicx}
\usepackage{asymptote}
\usepackage[hypertexnames=false]{hyperref} % make hyperref aware of chapter reset for built-in symbols
% (http://tex.stackexchange.com/questions/6098/wrong-hyper-references-after-resetting-chapter-counter)
\usepackage{listings}
\usepackage{paralist}
\usepackage{color}
\usepackage{fullpage}
\usepackage{graphics}
\usepackage{textcomp}
\usepackage{amsmath}
\usepackage{amssymb}
\usepackage{mathpazo}
\usepackage[mathpazo]{flexisym}
\usepackage{breqn}
\usepackage{environ}
\usepackage{multicol}
\usepackage[]{colophon}
\usepackage[k-tight]{minitoc}
\setlength{\mtcindent}{0pt}
\mtcsetformat{minitoc}{tocrightmargin}{2.55em plus 1fil}
\newcommand{\multicolumnmtc}{3}
\makeatletter
\let\SV@mtc@verse\mtc@verse
\let\SV@endmtc@verse\endmtc@verse
\def\mtc@verse#1{\SV@mtc@verse#1\removelastskip%
 \begin{multicols}{\multicolumnmtc}\raggedcolumns\leavevmode\unskip
 \vskip -1.5ex \vskip -1\baselineskip}
\def\endmtc@verse{\end{multicols}\SV@endmtc@verse}
\makeatother

\definecolor{message}{rgb}{0.8,0.4,0}
\definecolor{print}{rgb}{0,0,1}
\definecolor{definitions}{rgb}{0.96,0.96,0.94}
\definecolor{subtitle}{rgb}{0.3,0.3,0.3}

\definecolor{pythonstring}{rgb}{0,0.5,0}

\title{
    \includegraphics[height=0.1\linewidth]{logo-heptatom.pdf}
    \includegraphics[height=0.08125\linewidth]{logo-text-nodrop.pdf}
\\[.5em]
  {\LARGE\color{subtitle}\textit{\textmd{A free, open-source alternative to Mathematica}}}
  \par\textmd{\Large Mathics Core Version \MathicsCoreVersion}
}
\author{The Mathics Team}


% Fix unicode mappings for listings
% http://tex.stackexchange.com/questions/39640/typesetting-utf8-listings-with-german-umlaute
\lstset{literate=%
    {Ö}{{\"O}}1
    {Ä}{{\"A}}1
    {Ü}{{\"U}}1
    {ß}{{\ss}}1
    {ü}{{\"u}}1
    {ä}{{\"a}}1
    {ö}{{\"o}}1
    {~}{{\textasciitilde}}1
    {⧦}{{===}}1
    {}{{=>}}1
    {⇒}{{$\Rightarrow$}}1
    {⊻}{{$\oplus$}}1
    {…}{{...}}1
    {∫}{{S}}1
    {}{{d}}1
}

\setlength{\columnsep}{1cm}

\makeatletter\newenvironment{definitions}{

	\vspace{0.5em}%
	\begin{lrbox}{\@tempboxa}\begin{minipage}{0.95\columnwidth}%
	\nobreak\vspace{0.3em}%
}{%
	\nobreak%
	\vspace{0.3em}%
	\end{minipage}\end{lrbox}
	\colorbox{definitions}{\usebox{\@tempboxa}}
	\vspace{1em}%
}\makeatother

\newcommand{\dt}[1]{#1%
\par%
}
\newcommand{\dd}[1]{{%
\hangindent2em%
\hangafter=0%
#1%
\par%
}}

% Unfortunately, this doesn't work with code segments:
%\newenvironment{definitions}{\begin{description}}{\end{description}}
%\newcommand{\dt}[1]{\item[#1]}
%\newcommand{\dd}[1]{#1}

\newcommand{\sectionstart}{}
\newcommand{\sectionend}{}

\newcommand{\subsectionstart}{}
\newcommand{\subsectionend}{}

\newcommand{\chapterstart}{
}
\newcommand{\chaptersections}{
	\minitoc
	\begin{multicols}{2}
}
\newcommand{\chapterend}{
	\end{multicols}
}

\newcommand{\referencestart}{
\setcounter{chapter}{0}
%\def\thechapter{\Roman{chapter}}
\renewcommand{\chaptersections}{
	\minitoc
	\begin{multicols*}{2}
}
\renewcommand{\chapterend}{
	\end{multicols*}
}
}

% don't use two columns in appendix
\let\originalappendix\appendix
\renewcommand{\appendix}{\originalappendix
\renewcommand{\chaptersections}{
}
\renewcommand{\chapterend}{
}
}

\newenvironment{chapterintro}{\begin{multicols}{2}}{\end{multicols}}
\newenvironment{chapterintroshort}{\par}{}

\newenvironment{guidesectionintro}{\begin{multicols}{2}}{\end{multicols}}
\newenvironment{guidesectionintroshort}{\par}{}

\newcommand{\codewidth}{0.9\columnwidth}

\newenvironment{tests}{%
	\begin{compactitem}[{%
	\tiny\ttfamily $>>$
	}]%
}{%
	\end{compactitem}%
}
\newenvironment{testcase}{%
	\item%
	\begin{minipage}[t]{\codewidth}
	\begin{compactitem}[]%
}{%
	\end{compactitem}%
	\end{minipage}
	\vspace{1em}%
}

\newcommand{\test}[1]{%
\item \begin{minipage}[t]{\codewidth}\begin{flushleft} {\ttfamily \bfseries #1} \end{flushleft} \end{minipage} \hfill \nobreak%
}

\newenvironment{testmessage}{\item
\vspace{-0.5em}
\begin{minipage}{\codewidth}
\begin{dmath*}[indentstep={0em},compact]%
\color{message}%
}{%
\end{dmath*}%
\end{minipage}%
}
\newenvironment{testprint}{\item
\vspace{-0.5em}
\begin{minipage}{\codewidth}
\begin{dmath*}[indentstep={0em},compact]%
\color{print}%
}{%
\end{dmath*}%
\end{minipage}%
}

\newenvironment{testresult}{%
	\item%
	\vspace{-0.5em}%
	\begin{minipage}{\codewidth}%
	\hbadness=10000
	\begin{dmath*}[indentstep={0em},compact]%
}{%
	\end{dmath*}%
	\end{minipage}%
}
\newenvironment{testresultlist}{%
	\item%
	\vspace{-0.5em}%
	\begin{minipage}{\codewidth}%
	\begin{dgroup*}[indentstep={0em},compact]%
}{%
	\end{dgroup*}%
	\end{minipage}%
}

\newcommand{\code}[1]{\hbadness=10000{\ttfamily #1}}
\newcommand{\codevar}[1]{{\rmfamily #1}}
\newcommand{\console}[1]{\hbadness=10000{\ttfamily #1}}

\setlength{\parindent}{0mm}
\setlength{\parskip}{1pt}

\setlength{\mathindent}{0em}

\raggedbottom

\raggedcolumns

\makeindex

\hyphenation{assign-ments}

\begin{document}
\input{version-info.tex}

\maketitle

\dominitoc

\setcounter{tocdepth}{0}
\tableofcontents

\lstset{
%  inputencoding=utf8,
  extendedchars=true,
  breaklines=true,
  breakatwhitespace=false,
  showspaces=false,
  basicstyle=\ttfamily
}

\lstdefinestyle{python}{language=Python,tabsize=2,basicstyle=\scriptsize\ttfamily,showstringspaces=false,stringstyle=\color{pythonstring},commentstyle=\ttfamily\color{pythonstring},frame=lines}

% -*- latex -*-

% This file can be used for troubleshooting problems.

% Adjust this file to your liking and then build using:
%    make mathics-test.pdf
\chapter{Test Chapter}
\chapterstart
\chaptersections
\section*{TestSection}
\sectionstart

Here is what we want to test:

        \begin{asy}
        usepackage("amsmath");
        size(6.6667cm, 4.1202cm);
        // LineBox
        draw((18.182,231.407)--(25.603,235.593)--(33.024,235.315)--(40.445,230.584)--(47.866,221.587)--(55.288,208.685)--(62.709,192.39)--(70.13,173.354)--(77.551,152.334)--(84.972,130.169)--(92.393,107.742)--(99.814,85.9482)--(107.24,65.6553)--(114.66,47.6727)--(122.08,32.7175)--(129.5,21.3856)--(136.92,14.129)--(144.34,11.237)--(151.76,12.8248)--(159.18,18.8291)--(166.6,29.0105)--(174.03,42.9632)--(181.45,60.131)--(188.87,79.8293)--(196.29,101.273)--(203.71,123.607)--(211.13,145.941)--(218.55,167.384)--(225.97,187.083)--(233.4,204.25)--(240.82,218.203)--(248.24,228.385)--(255.66,234.389)--(263.08,235.977)--(270.5,233.085)--(277.92,225.828)--(285.34,214.496)--(292.76,199.541)--(300.19,181.558)--(307.61,161.265)--(315.03,139.471)--(322.45,117.045)--(329.87,94.8794)--(337.29,73.8596)--(344.71,54.8231)--(352.13,38.5288)--(359.55,25.6262)--(366.98,16.6298)--(374.4,11.8983)--(381.82,11.6202), rgb(0.2472, 0.24, 0.6)+linewidth(0.66667));
        // PolygonBox
        filldraw((18.182,231.41)--(25.603,235.59)--(33.024,235.32)--(40.445,230.58)--(47.866,221.59)--(55.288,208.68)--(62.709,192.39)--(70.13,173.35)--(77.551,152.33)--(84.972,130.17)--(92.393,107.74)--(99.814,85.948)--(107.24,65.655)--(114.66,47.673)--(122.08,32.717)--(129.5,21.386)--(136.92,14.129)--(144.34,11.237)--(151.76,12.825)--(159.18,18.829)--(166.6,29.011)--(174.03,42.963)--(181.45,60.131)--(188.87,79.829)--(196.29,101.27)--(203.71,123.61)--(211.13,145.94)--(218.55,167.38)--(225.97,187.08)--(233.4,204.25)--(240.82,218.2)--(248.24,228.38)--(255.66,234.39)--(263.08,235.98)--(270.5,233.08)--(277.92,225.83)--(285.34,214.5)--(292.76,199.54)--(300.19,181.56)--(307.61,161.27)--(315.03,139.47)--(322.45,117.04)--(329.87,94.879)--(337.29,73.86)--(344.71,54.823)--(352.13,38.529)--(359.55,25.626)--(366.98,16.63)--(374.4,11.898)--(381.82,11.62)--(381.82,0)--(18.182,0)--cycle, evenodd+rgb(0.2472, 0.24, 0.6)+opacity(0.2), nullpen);
        clip(box((-19.906,-13.611), (416.67,260.82)));
        // LineBox
        draw((-16,123.607)--(416,123.607), rgb(0, 0, 0)+linewidth(0.66667));
        // InsetBox
        label("$\text{10}$", (84.9721706864564,120.94013108331227), (0,-1), rgb(0, 0, 0));
        // InsetBox
        label("$\text{20}$", (159.18367346938774,120.94013108331227), (0,-1), rgb(0, 0, 0));
        // InsetBox
        label("$\text{30}$", (233.3951762523191,120.94013108331227), (0,-1), rgb(0, 0, 0));
        // InsetBox
        label("$\text{40}$", (307.6066790352504,120.94013108331227), (0,-1), rgb(0, 0, 0));
        // InsetBox
        label("$\text{50}$", (381.81818181818176,120.94013108331227), (0,-1), rgb(0, 0, 0));
        // LineBox
        draw((84.972,123.607)--(84.972,130.273), rgb(0, 0, 0)+linewidth(0.66667));draw((159.18,123.607)--(159.18,130.273), rgb(0, 0, 0)+linewidth(0.66667));draw((233.4,123.607)--(233.4,130.273), rgb(0, 0, 0)+linewidth(0.66667));draw((307.61,123.607)--(307.61,130.273), rgb(0, 0, 0)+linewidth(0.66667));draw((381.82,123.607)--(381.82,130.273), rgb(0, 0, 0)+linewidth(0.66667));draw((-4.0816,123.607)--(-4.0816,127.607), rgb(0, 0, 0)+linewidth(0.66667));draw((25.603,123.607)--(25.603,127.607), rgb(0, 0, 0)+linewidth(0.66667));draw((40.445,123.607)--(40.445,127.607), rgb(0, 0, 0)+linewidth(0.66667));draw((55.288,123.607)--(55.288,127.607), rgb(0, 0, 0)+linewidth(0.66667));draw((70.13,123.607)--(70.13,127.607), rgb(0, 0, 0)+linewidth(0.66667));draw((99.814,123.607)--(99.814,127.607), rgb(0, 0, 0)+linewidth(0.66667));draw((114.66,123.607)--(114.66,127.607), rgb(0, 0, 0)+linewidth(0.66667));draw((129.5,123.607)--(129.5,127.607), rgb(0, 0, 0)+linewidth(0.66667));draw((144.34,123.607)--(144.34,127.607), rgb(0, 0, 0)+linewidth(0.66667));draw((174.03,123.607)--(174.03,127.607), rgb(0, 0, 0)+linewidth(0.66667));draw((188.87,123.607)--(188.87,127.607), rgb(0, 0, 0)+linewidth(0.66667));draw((203.71,123.607)--(203.71,127.607), rgb(0, 0, 0)+linewidth(0.66667));draw((218.55,123.607)--(218.55,127.607), rgb(0, 0, 0)+linewidth(0.66667));draw((248.24,123.607)--(248.24,127.607), rgb(0, 0, 0)+linewidth(0.66667));draw((263.08,123.607)--(263.08,127.607), rgb(0, 0, 0)+linewidth(0.66667));draw((277.92,123.607)--(277.92,127.607), rgb(0, 0, 0)+linewidth(0.66667));draw((292.76,123.607)--(292.76,127.607), rgb(0, 0, 0)+linewidth(0.66667));draw((322.45,123.607)--(322.45,127.607), rgb(0, 0, 0)+linewidth(0.66667));draw((337.29,123.607)--(337.29,127.607), rgb(0, 0, 0)+linewidth(0.66667));draw((352.13,123.607)--(352.13,127.607), rgb(0, 0, 0)+linewidth(0.66667));draw((366.98,123.607)--(366.98,127.607), rgb(0, 0, 0)+linewidth(0.66667));draw((396.66,123.607)--(396.66,127.607), rgb(0, 0, 0)+linewidth(0.66667));
        // LineBox
        draw((10.761,-12.9443)--(10.761,260.158), rgb(0, 0, 0)+linewidth(0.66667));
        // InsetBox
        label("$\text{-1.0}$", (8.09400123685838,11.189047026645202), (-1,0), rgb(0, 0, 0));
        // InsetBox
        label("$\text{-0.5}$", (8.09400123685838,67.39792238831208), (-1,0), rgb(0, 0, 0));
        // InsetBox
        label("$\text{0.5}$", (8.09400123685838,179.81567311164582), (-1,0), rgb(0, 0, 0));
        // InsetBox
        label("$\text{1.0}$", (8.09400123685838,236.02454847331268), (-1,0), rgb(0, 0, 0));
        // LineBox
        draw((10.761,11.189)--(17.427,11.189), rgb(0, 0, 0)+linewidth(0.66667));draw((10.761,67.3979)--(17.427,67.3979), rgb(0, 0, 0)+linewidth(0.66667));draw((10.761,179.816)--(17.427,179.816), rgb(0, 0, 0)+linewidth(0.66667));draw((10.761,236.025)--(17.427,236.025), rgb(0, 0, 0)+linewidth(0.66667));draw((10.761,-0.052728)--(14.761,-0.052728), rgb(0, 0, 0)+linewidth(0.66667));draw((10.761,22.4308)--(14.761,22.4308), rgb(0, 0, 0)+linewidth(0.66667));draw((10.761,33.6726)--(14.761,33.6726), rgb(0, 0, 0)+linewidth(0.66667));draw((10.761,44.9144)--(14.761,44.9144), rgb(0, 0, 0)+linewidth(0.66667));draw((10.761,56.1561)--(14.761,56.1561), rgb(0, 0, 0)+linewidth(0.66667));draw((10.761,78.6397)--(14.761,78.6397), rgb(0, 0, 0)+linewidth(0.66667));draw((10.761,89.8815)--(14.761,89.8815), rgb(0, 0, 0)+linewidth(0.66667));draw((10.761,101.123)--(14.761,101.123), rgb(0, 0, 0)+linewidth(0.66667));draw((10.761,112.365)--(14.761,112.365), rgb(0, 0, 0)+linewidth(0.66667));draw((10.761,134.849)--(14.761,134.849), rgb(0, 0, 0)+linewidth(0.66667));draw((10.761,146.09)--(14.761,146.09), rgb(0, 0, 0)+linewidth(0.66667));draw((10.761,157.332)--(14.761,157.332), rgb(0, 0, 0)+linewidth(0.66667));draw((10.761,168.574)--(14.761,168.574), rgb(0, 0, 0)+linewidth(0.66667));draw((10.761,191.057)--(14.761,191.057), rgb(0, 0, 0)+linewidth(0.66667));draw((10.761,202.299)--(14.761,202.299), rgb(0, 0, 0)+linewidth(0.66667));draw((10.761,213.541)--(14.761,213.541), rgb(0, 0, 0)+linewidth(0.66667));draw((10.761,224.783)--(14.761,224.783), rgb(0, 0, 0)+linewidth(0.66667));draw((10.761,247.266)--(14.761,247.266), rgb(0, 0, 0)+linewidth(0.66667));
        \end{asy}

\sectionend
\chapterend


\printindex
\begin{colophon}
\begin{description}
  \item[Mathics Core] \hfill \\ \MathicsCoreVersion
  \item[Python] \hfill \\ \PythonVersion
  \item[XeTeX] \hfill \\ \XeTeXVersion
  \item[Asymptote] \hfill \\ \AsymptoteVersion
  \item[Ghostscript] \hfill \\ \GhostscriptVersion
\end{description}
\end{colophon}
\end{document}
